% Generated draft: requires mathematical and source review.
\documentclass[11pt]{article}
\usepackage[margin=25mm]{geometry}
\usepackage{fontspec}
\setmainfont{Latin Modern Roman}
\usepackage{amsmath,amssymb,mathtools,unicode-math}
\setmathfont{Latin Modern Math}
\usepackage{enumitem,needspace,xurl,hyperref}
\setlength{\parindent}{0pt}
\setlength{\parskip}{5pt}
\setlength{\emergencystretch}{3em}
\newcommand{\N}{\mathbb{N}}
\newcommand{\Z}{\mathbb{Z}}
\newcommand{\Q}{\mathbb{Q}}
\newcommand{\R}{\mathbb{R}}
\newcommand{\C}{\mathbb{C}}
\newcommand{\F}{\mathbb{F}}
\newcommand{\A}{\mathbb{A}}
\newcommand{\PP}{\mathbb{P}}
\newcommand{\E}{\mathbb{E}}
\newcommand{\eps}{\varepsilon}
\newcommand{\dd}{\,\mathrm{d}}
\begin{document}
\section*{UnsolvedMath ID: 11000066}
\textbf{Source problem code:} catalogue entry
\subsection{Definitions and mathematical statement}
Retain \(\mathfrak h_g\) as the Lie algebra called ``the Lie algebra for the Johnson filtration'' in the source. Its construction, coefficient ring or field, grading, and genus range are not included in the supplied excerpt. In particular, it is not safe to identify a Johnson-filtration associated graded algebra, its image under Johnson homomorphisms, the full symplectic-derivation algebra, and a Torelli Malcev Lie algebra without that definition.

\emph{Question 5.22.} Is \(\mathfrak h_g\) finitely presented as a Lie algebra? If so, give an explicit finite presentation. Over an intended field \(k\), finite presentation means existence of finite sets of generators \(x_1,\ldots,x_m\) and relations \(r_1,\ldots,r_s\) in their free Lie algebra such that
\[
 \mathfrak h_g\cong
 \mathcal L_k(x_1,\ldots,x_m)/\langle r_1,\ldots,r_s\rangle_{\mathrm{Lie}},
\]
where the denominator is the Lie ideal generated by the relations. The analogous coefficient-ring convention must be used if the source works integrally. If a graded presentation is intended, the grading of generators and homogeneous relations must also be specified. The excerpt does not require a single uniform finite presentation independent of \(g\); it asks finite presentation for the intended genus-dependent algebra. Finite-dimensionality of each graded piece, or finite generation alone, is not the requested finite-presentation conclusion.
\subsection{Short English statement}
Determine whether the source's Johnson-filtration Lie algebra \(\mathfrak h_g\) has finitely many generators and defining relations, and produce them if it does. Its exact definition and coefficient convention are missing, so it must not be replaced by another related Lie algebra.
\subsection{Sources}
\begin{itemize}
\item Local UnsolvedMath registry, permanent ID 11000066, source code catalogue entry, Question 5.22, ``Lie algebra for the Johnson filtration,'' PDF page 51. No individual proposer is supplied.
\item Farb, \emph{Problems on Mapping Class Groups and Related Topics}. Supplied URL: \url{https://www.math.uchicago.edu/~farb/papers/mcgbook.pdf}; recorded access 2026-07-29; extraction pdf-font-structure. The edited volume contains 329 explicitly typeset headings across chapters, with repeated local numbering disambiguated by page.
\item The supplied research summary refers to Hain's work, associated graded Johnson-filtration Lie algebras, and ``Andreadakis--Johnson'' questions, but provides no precise theorem, full reference, or definition of \(\mathfrak h_g\).
\item Supplied status open and OPEN-TRIAGE. The summary calls complete determination open, which is broader than this specific finite-presentation question and is not an independently verified current assessment. Status evidence and redistribution terms remain NEEDS\_REVIEW.
\end{itemize}
\end{document}
